\section{Data}
\label{sec:data}

\subsection{Six-task product evaluation (v2)}
\label{sec:evalv2}

We froze the evaluation before evaluating any model on it. Its sources are native Korean
text from KLUE machine reading comprehension (MRC) and Wizard of Seoul (WoS) dialogues
\citep{park2021klue}, drawn from components that no Bobcat training, calibration or
development data uses (6,028 rows), and a catalog of tool-call situations with explicit
(task, command) label tables (372 rows from 12 situations). There is no machine translation.
Table~\ref{tab:evalv2} lists the tasks.

\begin{table}[t]
\centering
\small
\caption{Six-task evaluation v2: 6,400 decisions in 2,328 components. ``Components'' are
contexts, dialogues or tool situations; splits are made at this level.}
\label{tab:evalv2}
\setlength{\tabcolsep}{4pt}
\begin{tabular}{@{}lrrp{4.0cm}p{4.0cm}@{}}
\toprule
Task & Rows & Comp. & Question families & Label basis \\
\midrule
Search passage selection & 1,726 & 712 & answers-query (Noul); title selection among 8, 77,
  200 or 255 candidates (Choice) & MRC answerable/unanswerable marks (human); contrast
  articles; source article \\
Citation verification & 1,100 & 400 & supports / contradicts / says nothing & MRC answers;
  altered numbers; other articles \\
Tool-call review & 372 & 12 & scope, irreversible, plan mismatch & catalog tables of
  (task, command) \\
External document screening & 1,110 & 222 & instruction to the system; evidence under
  injection & insertion rules for injected and benign notices; MRC answerability (human) \\
Request routing & 1,092 & 482 & service (5 services); reservation & human slot-domain
  annotation in WoS \\
Classification & 1,000 & 500 & fine taxonomy (9 labels); coarse taxonomy (5 labels) &
  publisher categories mapped to a new taxonomy \\
\midrule
Total & 6,400 & 2,328 & & \\
\bottomrule
\end{tabular}
\end{table}

\paragraph{Splits.}
Components are assigned to development (3,188 decisions, 1,161 components), calibration
(1,605) and final (1,607). No component crosses splits, there are no duplicate IDs, and rows
that use another article as a contrast use only articles from the same split. Tool-call
situations are assigned whole: six in development, three in final and the remaining three in
calibration, so the three final situations were never seen during development. Per-task
counts per split are in Appendix~\ref{app:tables}. Rebuilding with the recorded seed and
source files reproduces the SHA256 of every split file.

\paragraph{Counterfactual groups.}
2,193 groups change exactly one condition within a component: an answerable and an
unanswerable question on the same passage (human-written native pairs); a claim with the
original number and with an altered number; the same command under different tasks (for
example, a database reset command is \emph{expected} for a reset request and
\emph{overreach} for a request to add a column); and the same article under the fine and the
coarse taxonomy. We report the fraction of groups with a changing target in which every
member is answered correctly.

\paragraph{Invariance probes.}
The same query is asked with 8, 77, 200 and 255 title candidates, and the same evidence
judgment is asked before and after inserting injected text into the document.

\paragraph{Label reliability.}
Each row records how its label was obtained: 2,566 rows reuse the source's human
annotation, 3,534 rows follow a clear rule that can occasionally fail (a contrast article
may happen to contain the answer; publisher categories are noisy; catalog judgments), and
300 rows are number-altered \emph{contradicts} cases that become wrong if the question
admits several values. No row has been reviewed by a person. Reading samples confirmed that
the classification labels are noisy (for example, an article about a card company's
personnel reshuffle labelled as finance).

\paragraph{Leakage controls.}
We removed 27 MRC contexts whose normalized sentences exactly match KLUE NLI training
premises or hypotheses, or YNAT training titles of at least 15 characters. This does not
exclude near duplicates or exposure during the backbone's pretraining. Question and example
IDs, source GUIDs and answer strings never appear in model input; question IDs are fixed per
family. Rows were selected in an answer-independent hash order, except that routing uses a
per-service quota of 160 and external document screening is balanced between answerable and
unanswerable cases; both stratifications are recorded in the manifest.

\paragraph{What the evaluation does not cover.}
It has no Score primitive (the sources contain no native ordinal gold), no English side, and
no model-routing task (difficulty or required capability); WoS provides only service
routing. Title selection retains some ambiguity between articles on similar topics even
though contrast titles were filtered to share no words with the query.

\subsection{A lesson from version 1}
\label{sec:v1bug}

The first version of the evaluation stored rows with sorted JSON keys. This silently
reordered every Choice \texttt{criteria} map alphabetically, while the row's list of
candidate IDs kept the original order, so the order shown to the model and the gold index
disagreed. The build audit ran before the rows were written and could not see this. The bug
surfaced in the first readout (A.X-4.0-Light), which swapped two routing services almost
everywhere and scored 15.8\% on citation verification, below chance; its task macro on v1
was 43.1\%, against 62.4\% on v2. Noul rows, whose candidate order is fixed, were barely
affected. Version 2 preserves key order, stores a per-row hash that changes when the
candidate order changes, re-reads the saved files and compares the presented order with the
candidate IDs; the readout refuses v1 files and aborts on any order mismatch. The general
lesson is that audits must run on the exact bytes the model will read.

\subsection{Training mixture}
\label{sec:mixture}

\begin{table}[t]
\centering
\small
\caption{Training mixture of \expone{} (one epoch).}
\label{tab:mixture}
\begin{tabular}{@{}p{10.8cm}r@{}}
\toprule
Part & Rows \\
\midrule
Task-type rows from KLUE MRC and WoS components not used by the evaluation & 25,144 \\
Policy-transfer rows (prepared in an earlier stage of the project) & 7,752 \\
Public gold: NLI, sentiment, news topic, intent, QA, mean ratings (STS, HelpSteer) & 18,042 \\
\midrule
Total (31,034 components, 36.25M tokens) & 50,938 \\
\quad by primitive: Choice / Noul / Score & 25,621 / 20,082 / 5,235 \\
\quad by language: Korean / English & 41,621 / 9,317 \\
\midrule
Removed: task-type rows whose text overlaps evaluation text & 8,798 \\
Separate monitor split (training distribution) & 1,589 \\
\bottomrule
\end{tabular}
\end{table}

Table~\ref{tab:mixture} summarizes the mixture. The task-type rows are built by the same
generators as the evaluation, from components the evaluation does not use, and their
injected passages use wording different from the evaluation's. The public gold part adds
NLI, sentiment, news-topic, intent and QA decisions and mean ratings from STS and HelpSteer
\citep{wang2023helpsteer}; the mean ratings are supervised through the mean-rating loss of
Section~\ref{sec:objectives}. The classification task received only 326 rows.

\paragraph{Held-out task.}
The tool-call review task is excluded from training altogether, so its evaluation measures
transfer to an unseen task. The other five tasks share generators with the training data
(with disjoint components), so their gains are in-distribution.

\paragraph{No teacher or Jev signal.}
Although the project's plan designates GLM-5.3-Flash as a permitted teacher, \expone{} was
trained on gold labels only. Its manifest records that no teacher outputs and no Jev outputs
were used.
